\section{Step 14: Perform one reflected CRC bit recurrence}
Save the accumulator's low bit, then shift the accumulator right by one
with a logical 64-bit shift. If the saved bit was one, XOR in the reflected
polynomial. Decrement the bit counter. The supplied loop repeats until all eight
bits of the current input byte have been processed.

Fill \texttt{record\_crc\_bit}. The decision depends on the bit before the
shift. A conditional branch is a direct implementation. An equivalent version
can subtract that bit from zero to form either an all-zero or all-one mask, AND
it with the polynomial, and XOR the result after shifting. Compare work only
after both versions pass the same checksum and rejection tests.

\paragraph{Observe} Prepare \texttt{s14-1} and stop before
\texttt{record\_crc\_bit\_ready}. The first accumulator ends in a zero bit. At the first visit to
\texttt{record\_crc\_bit\_ready}, it is \texttt{0x000000007fffffff} and
\texttt{\$t4} is seven. Continue to observe all eight recurrences.

\Needspace{7\baselineskip}
\paragraph{Check} Which bit determines whether to XOR the polynomial?
\begin{itemize}
\item[A.] The old low bit.
\item[B.] The new high bit.
\item[C.] The next input byte's high bit.
\end{itemize}
